\chapter{大模型计算系统}
\label{ch:system}

大语言模型（\en{Large Language Model}，LLM）的参数规模可达数十亿乃至数千亿，训练和推理需要执行海量的张量运算，运行过程中还会不断产生激活值、梯度和缓存等中间状态。一个模型能否真正被训练出来、能否以可接受的速度和成本回答用户的问题，并不只取决于神经网络算法本身，还取决于承载它的整套计算系统：计算芯片负责执行矩阵运算，高带宽存储负责持续供给参数与缓存，互联网络负责在设备之间交换数据，软件框架负责把算法翻译成可执行的计算，推理服务和应用工具则把模型交付给最终用户。\term{大模型计算系统}就是由这些硬件与软件共同构成、相互配合的整体。

本章采用“需求—方案—效果”的思路展开。第\,\ref{sec:req}\,节首先回答“大模型对计算系统提出了哪些要求”，并给出几种可以手算的估算方法，由此引出计算系统的四层结构；第\,\ref{sec:hw}～\ref{sec:tools}\,节依次介绍各层的解决方案——算力资源、深度学习框架、推理服务和应用系统工具，并说明每种方案解决了什么问题、付出了什么代价；第\,\ref{sec:case}\,节通过一个完整的设计案例，把各层知识串联起来。

\begin{guide}
\guideitem{学习目标}{学完本章，读者应能够：}
\begin{enumerate}[leftmargin=2em]
  \item 用参数量、计算量和数据搬运量三个基本量，估算一个模型在训练或推理时对硬件的要求；
  \item 用算术强度和 Roofline 模型判断一项计算受算力限制还是受访存限制；
  \item 说明 CPU、GPU、NPU/ASIC 的分工，以及 HBM、NVLink、InfiniBand 分别作用于哪一段数据通路；
  \item 区分 CUDA 等计算软件栈、PyTorch 等深度学习框架与 vLLM、Ollama 等推理服务组件所处的层次；
  \item 估算训练显存与 KV 缓存，并据此选择分布式训练或推理优化手段；
  \item 根据延迟、吞吐、显存、成本和工程复杂度，为一个具体应用设计合适的系统方案。
\end{enumerate}
\end{guide}
